\section{Explicit prime-field seeds of every finite height}\label{fg:sec:seeds}
The seed required by Theorem~\ref{fg:thm:exact-seed} exists for every
prime and every finite height. We first construct an integer form whose
relevant diagonal coefficient is a single multinomial. The coefficient
preservation of Lemma~\ref{fg:lem:replica} then gives a smooth
anticanonical hypersurface over the prime field. We use the notation
$B_t$ of Section~\ref{fg:sec:counting}; the intermediate form need not
define a smooth variety, and we attach no formal group to it.

\subsection{An integer coefficient seed}
Fix a prime $p$ and an integer $h\ge2$, and put
\[
 N=p^h-1,\qquad s=(2h-1)(p-1).
\]
Let $\nu=(\nu_1,\ldots,\nu_s)$ be the nondecreasing list containing
\[
 \begin{array}{c|c}
 \text{entry}&\text{multiplicity}\\ \hline
 1&2p-1\\
 p^k,\quad 1\le k\le h-2&2p-2\\
 p^{h-1}&p-2 .
 \end{array}
\]
The middle range is empty when $h=2$, and the final multiplicity is
zero when $p=2$. Distinguish $\nu_1=1$, and write
\[
 w_j=\nu_{j+1}\quad(1\le j\le s-1),\qquad
 S_0=0,\quad S_j=w_1+\cdots+w_j.
\]
We have
\begin{equation}\label{fg:eq:seed-weight}
 \sum_{j=1}^s\nu_j
 =(2p-1)+2(p-1)\sum_{k=1}^{h-2}p^k+(p-2)p^{h-1}
 =N,\qquad S_{s-1}=N-1.
\end{equation}

For every integer $0\le t\le N-1$, define a binary vector
$u_t\in\{0,1\}^{s-1}$ as follows. If $t=S_j$, let $u_t$ have its first
$j$ entries equal to one and its remaining entries zero. At all other
indices, apply descending greedy selection to the list $w$: start with
residual $t$, visit the indices $s-1,s-2,\ldots,1$, and select $w_j$
exactly when it does not exceed the current residual, subtracting it
when selected. The selected indices are the one entries of $u_t$.
The proof below verifies that the residual ends at zero.

In $N$ variables define
\begin{equation}\label{fg:eq:seed-monomials}
 M_1=x_1^N,\qquad
 M_{j+1}=\prod_{t=1}^{N-1}
       x_{t+1}^{\,1+u_{t,j}-u_{t-1,j}}\quad(1\le j\le s-1),
 \qquad G_{p,h}=\sum_{j=1}^sM_j.
\end{equation}

\begin{proposition}[Exact diagonal coefficients of the seed]
\label{fg:prop:coefficient-seed}
The construction satisfies $w\cdot u_t=t$ for $0\le t\le N-1$.
The monomials in \eqref{fg:eq:seed-monomials} are pairwise distinct
of degree $N$, and every exponent in $M_2,\ldots,M_s$ is at most two.
Thus $G_{p,h}$ is a degree-$N$ integer form in $N$ variables, with
exactly $s$ terms, all with coefficient one.

For every integer $r\ge0$,
\begin{equation}\label{fg:eq:seed-period}
 B_{r+1}(G_{p,h})=
 \begin{cases}
 0,&N\nmid r,\\[1mm]
 \displaystyle\frac{(aN)!}{\prod_{j=1}^s(a\nu_j)!},
       &r=aN,\quad a\ge0.
 \end{cases}
\end{equation}
In particular,
\begin{equation}\label{fg:eq:seed-height-data}
 B_{p^i}(G_{p,h})=0\quad(1\le i<h),\qquad
 v_p\bigl(B_{p^h}(G_{p,h})\bigr)=h-1.
\end{equation}
The form has a division-free circuit with
$O(Ns+\log N)$ addition and multiplication gates, and its full exponent
table and such a circuit can be generated in time polynomial in
$N,s,\log N$.
\end{proposition}

\begin{proof}
The length of the prescribed list is
\[
 (2p-1)+(h-2)(2p-2)+(p-2)=(2h-1)(p-1)=s,
\]
and summing its entries gives \eqref{fg:eq:seed-weight}.
The remaining list $w$ has $2p-2$ unit entries. Before the first entry
of weight $p^k$, its lower-weight entries have total
\[
 2(p-1)(1+p+\cdots+p^{k-1})=2(p^k-1).
\]
Consequently, each entry of this sorted list is at most one plus the
sum of all preceding entries.

For any sorted positive integer list satisfying this condition and
starting with one, descending greedy selection represents every
integer between zero and the total weight. Indeed, if a target is at
least the final entry, selecting that entry leaves a nonnegative
target at most the preceding total. If it is smaller than the final
entry, it is already at most the preceding total, by the displayed
condition. Induction proves the assertion. The prescribed prefix
vectors also have their required weighted sums, and the positive
weights make all prefix indices distinct. Thus
\[
 w\cdot u_t=t,\qquad u_0=0,\qquad
 u_{N-1}=(1,\ldots,1).
\]

Each exponent in $M_{j+1}$ lies in $\{0,1,2\}$, and telescoping gives
\[
 \deg M_{j+1}
 =\sum_{t=1}^{N-1}(1+u_{t,j}-u_{t-1,j})
 =N-1+u_{N-1,j}-u_{0,j}=N.
\]
Only $M_1$ involves $x_1$. If two other exponent columns, indexed by
$j<k$, were equal, their successive differences would give
$u_{t,j}=u_{t,k}$ for every $t$, starting from $u_0=0$. At $t=S_j$,
however, the prescribed prefix vector has $u_{t,j}=1$ and $u_{t,k}=0$.
This proves pairwise distinctness and the coefficient-one assertion.

Consider a term contributing to
$(x_1\cdots x_N)^r$ in $G_{p,h}^r$, and let $n_j$ be the number of
selected copies of $M_j$. Its $x_1$ exponent is $Nn_1$. Therefore
there is no contribution unless $r=aN$ for an integer $a\ge0$.
In that case
\[
 n_1=a,\qquad \sum_{j=1}^{s-1}n_{j+1}=a(N-1).
\]
Write $n'=(n_2,\ldots,n_s)$. The required exponent $aN$ at
$x_{t+1}$ gives
\[
 a(N-1)+(u_t-u_{t-1})\cdot n'=aN,
 \qquad (u_t-u_{t-1})\cdot n'=a.
\]
Summing from $u_0=0$ yields $u_t\cdot n'=at$. At the prefix indices
$t=S_j$, this is
\[
 n_2+\cdots+n_{j+1}=aS_j.
\]
Successive differences force $n_{j+1}=aw_j$, so the only possible
count vector is $n=a\nu$. Conversely, this vector has the required
total and all the required exponents because $w\cdot u_t=t$.
The multinomial theorem therefore proves the exact integer identity
\eqref{fg:eq:seed-period}, including its $a=0$ case.

For $1\le i<h$, we have $0<p^i-1<N$, so
$B_{p^i}(G_{p,h})=0$. For $i=h$, the coefficient is
$N!/\prod_j\nu_j!$. Legendre's formula, already used in
Lemma~\ref{fg:lem:replica}, gives
\[
 v_p\left(\frac{N!}{\prod_j\nu_j!}\right)
 =\frac{\sum_j s_p(\nu_j)-s_p(N)}{p-1}
 =\frac{(2h-1)(p-1)-h(p-1)}{p-1}=h-1.
\]
Here each $\nu_j$ is a power of $p$, and the $h$ base-$p$ digits of
$N=p^h-1$ are all $p-1$. This proves
\eqref{fg:eq:seed-height-data}.

Finally, there are $N$ vectors $u_t$, each obtained by at most
$s-1$ comparisons and subtractions on $O(\log N)$-bit integers.
The exponent table follows by adjacent differences. Precompute the
squares $x_i^2$; each of the last $s-1$ monomials then uses $O(N)$
multiplications, while $x_1^N$ uses $O(\log N)$ by repeated squaring.
Adding the terms proves the gate and construction bounds.
\end{proof}

For instance, $p=2,h=2$ gives
$G_{2,2}=x_1^3+x_2^2x_3+x_2x_3^2$, with $B_2=0$ and $B_4=6$.
At $p=2,h=3$, the multiplicity list is $(1,1,1,2,2)$, and
$B_8=7!/(2!2!)=1260$ has valuation two. These values illustrate the
integer identity; no smoothness claim about an intermediate form is
needed.

\subsection{Smooth seeds and exact-height completeness}
\begin{theorem}[Every finite height over the prime field]
\label{fg:thm:all-heights}
For every prime $p$ and every finite integer $h\ge1$, there is an
explicit smooth geometrically integral anticanonical projective
hypersurface over $\Fp$ whose top Artin--Mazur formal group has exact
height $h$. One may take its degree and number of homogeneous
variables to be
\begin{equation}\label{fg:eq:smooth-seed-size}
 d_{p,h}=p^h+p+1.
\end{equation}
Its dimension is then $p^h+p-1$.
\end{theorem}

\begin{proof}
First assume $h\ge2$ and take the form $G_{p,h}$ above, in
$N=p^h-1$ variables. Choose
\[
 Q=p^h+p.
\]
This is a positive multiple of $p$ with
$Q\ge N+2$ and $Q\ge p^h$. Theorem~\ref{thm:stabilizer} applies to
the reduction of $G_{p,h}$ over $\Fp$ and gives a smooth
geometrically integral hypersurface $X=V(F_{G_{p,h}})$ of degree
$Q+1$ in $Q+1$ variables. Thus it is anticanonical and has the
degree and dimension in \eqref{fg:eq:smooth-seed-size}.
The theorem allows every positive multiple $Q$ of $p$ satisfying
the size condition; $Q$ need not be a power of $p$.

Use the displayed integer stabilizer equation as the homogeneous
lift. Lemma~\ref{fg:lem:replica} gives exact equalities over the
integers
\[
 b_{p^i}(F_{G_{p,h}})=B_{p^i}(G_{p,h})
 \qquad(1\le i\le h).
\]
The earlier numerators vanish and the last has valuation $h-1$ by
\eqref{fg:eq:seed-height-data}. The integral Stienstra logarithm
and \eqref{fg:eq:exact-height} therefore give $h(X)=h$.
For $h=1$, use instead the degree-one form $G=x_1$ in one
variable and choose $Q=2p$. Then $Q\ge3$ and $Q\ge p$, so the
same stabilization and coefficient identity apply. Since
$B_p(x_1)=1$, the resulting smooth hypersurface has $b_p=1$ and
height one by Proposition~\ref{fg:prop:height-one}. Its degree and
variable count are $Q+1=2p+1=p^1+p+1$ as well. The formal-group
interpretation is used only after stabilization to positive dimension.

The construction is explicit by Proposition~\ref{fg:prop:coefficient-seed}
and the circuit construction in Theorem~\ref{thm:stabilizer}.
For each fixed $p,h$, its equation is fixed finite data.
\end{proof}

\begin{corollary}[Completeness at every fixed exact height]
\label{fg:cor:exact-height-complete}
Fix a prime $p$ and a finite integer $h\ge1$.
For division-free arithmetic circuits over $\Fp$ promised to define
smooth geometrically integral anticanonical hypersurfaces of positive
dimension, with dimension allowed to grow, deciding $h(X)=h$ is
$\ModP$-complete under deterministic polynomial-time many-one
reductions. Hardness persists after imposing any fixed lower bound
on the dimension.
\end{corollary}

\begin{proof}
Theorem~\ref{fg:thm:all-heights} supplies the fixed seed in
Theorem~\ref{fg:thm:exact-seed}, which proves completeness.
In its reduction, the smoothing parameter can also exceed any
prescribed fixed dimension bound. This preserves all coefficient
identities through $p^h$ and the polynomial circuit-size bound.
\end{proof}

The degree in \eqref{fg:eq:smooth-seed-size} is exponential in $h$.
The completeness statement fixes $p$ and $h$; it does not assert a
polynomial-size seed as a function of a binary-encoded varying height.

For the literature comparison,
\citet[Convention~4.15 and Theorem~4.18, pp.~25--28]{fg-kty-two}
construct anticanonical hypersurfaces of every finite height over an
algebraically closed field, using degree $p^h-1$ for $h>1$.
That proof intersects nonempty smoothness and height loci in a
parameter space. Its stated field convention does not supply a
prime-field point in that intersection. The coefficient identity
\eqref{fg:eq:seed-period} and the prescribed stabilization above
give equations over $\Fp$, with the explicit larger degree
\eqref{fg:eq:smooth-seed-size}.
